% ============================================================================
% End-Semester Review Presentation (Beamer 16:9, "clean card" theme: white, teal-green accent).
% Numbers identical to main_v2.tex (verified against source files).
% Build: pdflatex presentation_final.tex (twice)
% ============================================================================
\documentclass[aspectratio=169,10pt]{beamer}
\usepackage[utf8]{inputenc}
\usepackage[T1]{fontenc}
\usepackage[default]{sourcesanspro}
\usepackage{booktabs,array,tabularx,multicol}
\usepackage{graphicx}
\usepackage[most]{tcolorbox}
\graphicspath{{./}{figures/}}

\definecolor{accent}{HTML}{0F766E}   % deep teal-green
\definecolor{accent2}{HTML}{B45309}  % amber
\definecolor{keygreen}{HTML}{15803D} % key-finding stripe
\definecolor{ink}{HTML}{1F2937}
\definecolor{muted}{HTML}{6B7280}
\definecolor{card}{HTML}{F5F8F7}
\definecolor{cardwarm}{HTML}{FDF8F1}
\definecolor{rule}{HTML}{D1D5DB}
\colorlet{navy}{accent}
\colorlet{good}{keygreen}

\setbeamercolor{normal text}{fg=ink,bg=white}
\setbeamercolor{structure}{fg=accent}
\setbeamercolor{alerted text}{fg=accent2}
\setbeamercolor{footbar}{fg=muted,bg=white}
\setbeamertemplate{navigation symbols}{}
\setbeamertemplate{itemize item}{\scriptsize\color{accent}$\blacktriangleright$}
\setbeamertemplate{itemize subitem}{\color{accent}--}
\setbeamertemplate{enumerate item}{\color{accent}\bfseries\insertenumlabel.}
\setlength{\leftmargini}{1.4em}

% title: dark text, thin teal rule, small logo at right
\setbeamertemplate{frametitle}{%
  \vspace{3mm}%
  \begin{minipage}[c]{0.88\textwidth}{\Large\bfseries\color{ink}\insertframetitle}\end{minipage}\hfill
  \begin{minipage}[c]{0.09\textwidth}\raggedleft\includegraphics[height=7.5mm]{figures/logo.jpeg}\end{minipage}\par
  \vspace{0.5mm}{\color{accent}\rule{\textwidth}{1.4pt}}}

% footer: thin grey line, muted text, slide number
\setbeamertemplate{footline}{%
  \begin{beamercolorbox}[wd=\paperwidth,ht=4.4ex,dp=1.6ex,leftskip=0.45cm,rightskip=0.45cm]{footbar}%
    \vbox{\hbox to \dimexpr\paperwidth-0.9cm\relax{\color{rule}\leaders\hrule height 0.6pt\hfill}\vskip 1.2ex
    \hbox to \dimexpr\paperwidth-0.9cm\relax{\scriptsize IIITDM Kurnool \;\textbar\; Multimodal Agricultural Drought Monitoring and Prediction\hfill\insertframenumber\,/\,\inserttotalframenumber}}%
  \end{beamercolorbox}}

% cards: light background + coloured stripe on the left + coloured heading
\newtcolorbox{card}[2][]{enhanced, colback=card, colframe=card, boxrule=0pt, arc=2pt,
  borderline west={2.5pt}{0pt}{accent}, left=7pt, right=6pt, top=3pt, bottom=4pt,
  title={#2}, coltitle=accent, fonttitle=\bfseries\small, attach title to upper={\par\vspace{1pt}},
  before skip=3pt, after skip=5pt, #1}
\newtcolorbox{warncard}[2][]{enhanced, colback=cardwarm, colframe=cardwarm, boxrule=0pt, arc=2pt,
  borderline west={2.5pt}{0pt}{accent2}, left=7pt, right=6pt, top=3pt, bottom=4pt,
  title={#2}, coltitle=accent2, fonttitle=\bfseries\small, attach title to upper={\par\vspace{1pt}},
  before skip=3pt, after skip=5pt, #1}
\newtcolorbox{keycard}[2][]{enhanced, colback=card, colframe=card, boxrule=0pt, arc=2pt,
  borderline west={2.5pt}{0pt}{keygreen}, left=7pt, right=6pt, top=3pt, bottom=4pt,
  title={#2}, coltitle=keygreen, fonttitle=\bfseries\small, attach title to upper={\par\vspace{1pt}},
  before skip=3pt, after skip=5pt, #1}

\newcolumntype{Y}{>{\raggedright\arraybackslash}X}
\newcommand{\key}[2]{{\color{accent}\bfseries\Large #1}\\[-1pt]{\scriptsize\color{muted} #2}}
\setlength{\tabcolsep}{4pt}

\begin{document}

% ============================================================ 1 TITLE
\begin{frame}[plain]
\centering
\vspace{1mm}
\includegraphics[height=1.2cm]{figures/logo.jpeg}\\[1.5mm]
{\footnotesize\bfseries\color{muted} INDIAN INSTITUTE OF INFORMATION TECHNOLOGY DESIGN AND MANUFACTURING KURNOOL}\\[1mm]
{\color{accent}\rule{0.9\textwidth}{1.2pt}}\\[3mm]
\parbox{0.9\textwidth}{\centering\fontsize{13}{16.5}\selectfont\bfseries\color{ink} A Multimodal Framework for Agricultural Drought\\Monitoring and Prediction using Weather Time Series,\\Satellite Imagery, and UAV Semantic Segmentation\par}\\[2.5mm]
{\normalsize\color{accent}\bfseries End-Semester Review Presentation}\\[1.5mm]
{\color{accent}\rule{0.9\textwidth}{1.2pt}}\\[3mm]
{\footnotesize\color{muted} Presented by}\\[0.5mm]
{\small\textbf{G. Aryan Rajesh} (123CS0020) \quad \textbf{M. Rovan} (123CS0036) \quad \textbf{M. Harsha Vardhan} (123CS0044)}\\[3mm]
{\footnotesize\color{muted} Under the guidance of}\\[0.5mm]
{\small\bfseries Dr. M. Naresh Babu}, {\small Associate Professor, Dept.\ of CSE}\\[3mm]
{\footnotesize\color{muted} Department of Computer Science and Engineering \quad$\cdot$\quad November 2026}
\end{frame}

% ============================================================ 3 INTRODUCTION
\begin{frame}{1. Introduction: Agricultural Drought in Rayalaseema}
\vspace{-2.5mm}
\begin{columns}[T,onlytextwidth]
\column{0.455\textwidth}
\begin{card}[equal height group=intro]{The regional problem}
\footnotesize
\begin{itemize}\setlength{\itemsep}{2pt}
\item Rayalaseema (Kurnool, Anantapur, Sri Sathya Sai) and Sira (Karnataka): low, erratic rainfall.
\item \textbf{Kharif} (Jun--Nov) crops depend on the monsoon; \textbf{rabi} (Oct--Apr) crops on stored soil water.
\item Drought is declared by Government Orders (G.O.s) \alert{\textbf{late in the season}}; detection is weakest in AP and Karnataka.
\item Weather, satellite and drone studies are usually done \textbf{separately}.
\end{itemize}
\end{card}
\column{0.525\textwidth}
\begin{card}[equal height group=intro]{8 study sites}
\footnotesize
\renewcommand{\arraystretch}{1.2}\setlength{\tabcolsep}{5pt}
\begin{tabular*}{\linewidth}{@{}l@{\extracolsep{\fill}}lr@{}}
\textbf{Somandepalle} & {\color{muted}Sri Sathya Sai} & {\scriptsize 13.999N, 77.625E} \\
\textbf{Lepakshi} & {\color{muted}Sri Sathya Sai} & {\scriptsize 13.800N, 77.619E} \\
\textbf{Cholasamudram} & {\color{muted}Sri Sathya Sai} & {\scriptsize 13.812N, 77.569E} \\
\textbf{Sira} & {\color{muted}Tumakuru, KA} & {\scriptsize 13.798N, 76.975E} \\
\textbf{Peddakadubur} & {\color{muted}Kurnool} & {\scriptsize 15.783N, 77.253E} \\
\textbf{Dinnedevarapadu} & {\color{muted}Kurnool} & {\scriptsize 15.772N, 78.047E} \\
\textbf{Near Balaji Villas} & {\color{muted}Kurnool} & {\scriptsize 15.783N, 78.056E} \\
\textbf{Near Pullaiah College} & {\color{muted}Kurnool} & {\scriptsize 15.798N, 78.081E} \\
\end{tabular*}\\[3pt]
{\scriptsize\color{muted} Each site is a UAV survey field; 6 independent weather cells.}
\end{card}
\end{columns}
\begin{card}{Our idea}
\footnotesize \textbf{Weather predicts} where to look $\rightarrow$ \textbf{satellite confirms} which fields are stressed $\rightarrow$ \textbf{drone measures} the damage --- in \textbf{both seasons} (kharif and rabi).
\end{card}
\end{frame}

% ============================================================ 4 GAPS
\begin{frame}[t]{2. Research Gaps and Our Solutions}
\vspace{-5mm}
\begin{columns}[T,onlytextwidth]
\column{0.49\textwidth}
\begin{warncard}[equal height group=gaprow1]{Gap 1 -- Weak forecast evaluation}
\small\raggedright\hyphenpenalty=10000 Drought indices are often calibrated on the full record (test years leak into labels), and models are compared only with other ML models.\par\vspace{3pt}
\textcolor{good}{\textbf{Ours:}} SPEI frozen on 1979--2011; baselines include persistence and a physical water-balance model.
\end{warncard}
\vspace{1mm}
\begin{warncard}[equal height group=gaprow2]{Gap 2 -- Index-vs-index validation}
\small\raggedright\hyphenpenalty=10000 Satellite studies validate one index against another at coarse scale.\par\vspace{3pt}
\textcolor{good}{\textbf{Ours:}} 240 random 10\,m fields; replay against 4 seasons of government orders with real data delays.
\end{warncard}
\column{0.49\textwidth}
\begin{warncard}[equal height group=gaprow1]{Gap 3 -- Rabi drought unchecked}
\small\raggedright\hyphenpenalty=10000 Rabi drought is assessed after the season, with rain indices rarely checked against crop outcomes.\par\vspace{3pt}
\textcolor{good}{\textbf{Ours:}} soil-moisture deficit indicator validated against district crop yields (1997--2022).
\end{warncard}
\vspace{1mm}
\begin{warncard}[equal height group=gaprow2]{Gap 4 -- UAV not scalable, modules not linked}
\small\raggedright\hyphenpenalty=10000 UAV studies classify small images, do not report areas, and are not linked to weather or satellite.\par\vspace{3pt}
\textcolor{good}{\textbf{Ours:}} 7-class segmentation of orthomosaics up to 2.9 billion pixels, with m$^2$ areas; kharif and rabi hand fields to the drone.
\end{warncard}
\end{columns}
\end{frame}

% ============================================================ 5 PROBLEM / OBJECTIVES
\begin{frame}{3. Problem Statement and Objectives}
\begin{card}{Problem statement}
\footnotesize Using daily weather (1979--2026), satellite imagery (2000--2026) and drone imagery for 8 agricultural sites in Kurnool, Sri Sathya Sai (Andhra Pradesh) and Sira (Karnataka), build and honestly evaluate a framework that \textbf{forecasts} drought, \textbf{confirms} it on crops and \textbf{measures} it on fields, in both the kharif and rabi seasons.
\end{card}
\begin{columns}[T,onlytextwidth]
\column{0.49\textwidth}
\begin{card}{O1 -- Weather forecast}
\footnotesize Forecast the SPEI-30 drought class (healthy / moderate / severe) 7, 15 and 30 days ahead and beat persistence and a physical water-balance model.
\end{card}
\begin{card}{O2 -- Satellite confirmation}
\footnotesize Confirm crop stress on randomly sampled Sentinel-2 fields and replay the chain against government drought orders with real data delays.
\end{card}
\column{0.49\textwidth}
\begin{card}{O3 -- Rabi drought}
\footnotesize Build a rabi drought indicator from root-zone soil moisture and validate it against district crop yields.
\end{card}
\begin{card}{O4 -- UAV segmentation}
\footnotesize Segment drone orthomosaics into 7 classes (soil, crop, cotton boll, green weed, dry weed, drought, background) and estimate class areas in m$^2$.
\end{card}
\begin{card}{O5 -- Integration}
\footnotesize Connect the modules: kharif and rabi outputs list the fields to survey with the drone.
\end{card}
\end{columns}
\end{frame}

% ============================================================ 6 DATA + ARCHITECTURE
\begin{frame}{4. Data Sources}
\begin{card}{Data (all public or collected; nothing simulated)}
\small
\renewcommand{\arraystretch}{1.3}
\begin{tabularx}{\linewidth}{@{}p{4.2cm}Y@{\hspace{6pt}}r@{}}
\toprule
\textbf{Source} & \textbf{Content} & \textbf{Period} \\ \midrule
AgERA5 & daily rain, ET$_0$, temperature, humidity & 1979--2026 \\
Sentinel-2 (Earth Engine) & 10\,m NDVI: 240 fields, 1,350 points & 2019--2026 \\
ERA5-Land & root-zone soil moisture, 45 points & 1990--2023 \\
MODIS MOD13Q1 & 250\,m NDVI on cropland & 2000--2026 \\
District crop statistics & rabi yield of all crops & 1997--2022 \\
Drought G.O.s & declared mandals (comparison only) & 2023--2026 \\
UAV surveys (DJI Mini 4 Pro) & 2,254 photos, 10 locations (cotton, maize) $\rightarrow$ 14 orthomosaics\newline {\color{muted}flown by Dr.\ M.\ Naresh Babu \& team} & Nov 2025 -- Jan 2026 \\
\bottomrule
\end{tabularx}
\end{card}
\end{frame}

% ============================================================ 6b ARCHITECTURE
\begin{frame}{4. Overall Architecture}
\centering
\vspace{-1mm}
\includegraphics[width=\textwidth]{figures/v2_system_slide.png}
\end{frame}

% ============================================================ 7 LIT WEATHER
\begin{frame}{5.1 Literature Review -- Weather Drought Forecasting}
\vspace{1mm}
\footnotesize
\renewcommand{\arraystretch}{1.45}\setlength{\tabcolsep}{5pt}
\begin{tabularx}{\textwidth}{|>{\raggedright\arraybackslash}p{3.2cm}|Y|Y|Y|}
\hline
\textbf{Paper} & \textbf{Contribution} & \textbf{Limitation} & \textbf{Our response} \\ \hline
\textbf{SPEI} \newline {\color{muted}Vicente-Serrano et al., 2010} & Water-balance drought index (rain -- PET) & Fitted on the full record, so test years leak into labels & Fitted on 1979--2011 only, then frozen \\ \hline
\textbf{LSTM SPEI forecast} \newline {\color{muted}Dikshit et al., 2021} & LSTM forecasts monthly SPEI & Monthly only; no data delay & Daily forecast with an 8-day data delay \\ \hline
\textbf{Deep drought prediction} \newline {\color{muted}Marusov et al., 2024} & Deep networks on gridded weather & Compared mostly with ML models & Physics + ML hybrid; their protocol as a check \\ \hline
\textbf{Wavelet ANN / SVR} \newline {\color{muted}Belayneh et al., 2014} & Wavelet ANN and SVR forecast SPI & Rain-only index; no physical baseline & Persistence, Markov, climatology and water-balance baselines \\ \hline
\end{tabularx}
\vspace{3mm}
\begin{card}{Our direction}
\footnotesize 30-day window of daily weather + SPEI history + \textbf{water-balance physics} $\rightarrow$ drought class 7 / 15 / 30 days ahead, with leakage-free chronological testing.
\end{card}
\end{frame}

% ============================================================ 8 WEATHER METHOD
\begin{frame}{5.2 Weather Module -- Method}
\vspace{1mm}
\begin{columns}[T,onlytextwidth]
\column{0.49\textwidth}
\begin{card}[equal height group=wm1]{Labels: SPEI-30 drought class}
\small\raggedright
Daily water balance: $D_t = P_t - \mathrm{ET}_{0,t}$\\[3pt]
30-day sum: $\mathrm{WB30}_t = \sum_{k=0}^{29} D_{t-k}$\\[3pt]
$\mathrm{SPEI} = \Phi^{-1}\!\big(F(\mathrm{WB30})\big)$, log-logistic $F$ fitted per site and month on \textbf{1979--2011} only, then frozen.\\[4pt]
\mbox{\textbf{Healthy} $\ge 0$} \;$\cdot$\; \mbox{\textbf{Moderate} $-1.5$ to $0$}\\ \mbox{\textbf{Severe} $< -1.5$}
\end{card}
\begin{card}[equal height group=wm2]{Data split (purged)}
\small\raggedright
Train 1979--2011 \;$\cdot$\; validation 2012--2018\\ Test 2019--2026; no sample crosses a split.
\end{card}
\column{0.49\textwidth}
\begin{card}[equal height group=wm1]{Inputs and models}
\small
\begin{itemize}
\item 30-day window: 9 weather variables, SPEI history, day of year
\item \textbf{Water-balance ensemble}: fill the unknown future days with each 1979--2011 year $\rightarrow$ class probabilities
\item RF, XGBoost, LSTM, Transformer (5 seeds) $\rightarrow$ soft voting
\end{itemize}
\end{card}
\begin{card}[equal height group=wm2]{Baselines}
\small
Persistence, Markov, climatology and the water-balance ensemble (physics).
\end{card}
\end{columns}
\end{frame}

% ============================================================ 8b WEATHER ARCHITECTURE
\begin{frame}{5.3 Weather Module -- Architecture}
\centering
\vspace{2mm}
\includegraphics[width=\textwidth]{figures/v2_weather_slide.png}
\end{frame}

% ============================================================ 9 WEATHER RESULTS
\begin{frame}{5.4 Weather Module -- Results (test years 2019--2026)}
\begin{columns}[T,onlytextwidth]
\column{0.56\textwidth}
\scriptsize
\renewcommand{\arraystretch}{1.0}
\resizebox{0.9\linewidth}{!}{%
\begin{tabular}{@{}lccc@{}}
\toprule
\textbf{Macro-F1} & \textbf{7 days} & \textbf{15 days} & \textbf{30 days} \\ \midrule
Persistence (class today) & 0.662 & 0.515 & 0.424 \\
Markov transition & 0.493 & 0.424 & 0.377 \\
Climatology & 0.309 & 0.310 & 0.310 \\
Water-balance ensemble (physics) & 0.783 & 0.592 & 0.375 \\ \midrule
Random Forest & 0.775 & 0.617 & 0.383 \\
XGBoost & 0.814 & 0.631 & 0.418 \\
LSTM & 0.806 & 0.642 & 0.406 \\
Transformer & 0.809 & 0.650 & 0.407 \\
\textbf{Soft-voting ensemble (hybrid)} & \textbf{0.819} & \textbf{0.651} & \textbf{0.424} \\ \midrule
Balanced accuracy / ROC-AUC & 85.6\% / 0.97 & 68.3\% / 0.90 & 42.5\% / 0.68 \\
\bottomrule
\end{tabular}}\\[2pt]
\centering\includegraphics[width=\linewidth]{figures/res_w_bars.png}
\column{0.41\textwidth}
\begin{keycard}{Key findings}
\footnotesize
\begin{itemize}
\item Hybrid \textbf{0.819} vs persistence 0.662 at 7 days; \textbf{0.651} vs 0.515 at 15 days.
\item Beats the physics model significantly at 15 days (95\% CI [+0.010, +0.099]).
\item Water-balance probabilities are the most valuable input (+0.063 at 7 days).
\item Severe drought: \textbf{80\%} of severe days caught 7 days ahead (false-alarm ratio 0.34).
\end{itemize}
\end{keycard}
\end{columns}
\end{frame}

% ============================================================ 10 LIT SATELLITE
\begin{frame}{5.5 Literature Review -- Satellite Drought Monitoring}
\footnotesize
\renewcommand{\arraystretch}{1.22}
\begin{tabularx}{\textwidth}{|>{\raggedright\arraybackslash}p{3.3cm}|Y|Y|Y|}
\hline
\textbf{Paper} & \textbf{Contribution} & \textbf{Limitation} & \textbf{Our response} \\ \hline
\textbf{VCI} \newline Kogan, 1995 & NDVI scaled between historical minimum and maximum & Needs long records; with 4--5 years many dates get VCI $\approx 0$ by chance & VCI only as input; SVI as label \\ \hline
\textbf{SVI} \newline Peters et al., 2002 & z-score of NDVI for each date & Usually at coarse scale & Per field and day of year, scored year removed from the normal \\ \hline
\textbf{ML index blending} \newline Park et al., 2016 & RF / boosting on 16 MODIS + TRMM drought factors & Coarse; index-vs-index validation & 10\,m fields validated against government orders \\ \hline
\textbf{Sentinel-2 + ML, India} \newline Sharma et al., 2025 & Sentinel-2 indices with RF / XGBoost for drought detection & District ground truth; no forecasting & Chronological forecast and replay with real data delays \\ \hline
\end{tabularx}
\begin{card}{Our direction}
\footnotesize 240 randomly sampled fields $\rightarrow$ Sentinel-2 SVI $\rightarrow$ lag to weather, crop-condition forecast, and replay against drought G.O.s.
\end{card}
\end{frame}

% ============================================================ 11 SATELLITE
\begin{frame}{5.6 Satellite Module -- Architecture}
\centering
\includegraphics[height=0.76\textheight]{figures/v2_satellite_slide.png}
\end{frame}

% ============================================================ 12b SATELLITE METHOD + RESULTS
\begin{frame}{5.7 Satellite Module -- Method and Results}
\vspace{1mm}
\begin{columns}[T,onlytextwidth]
\column{0.47\textwidth}
\begin{card}[equal height group=sat]{Fields and label}
\small
\begin{itemize}
\item Farmland within 3\,km of each site mapped in Earth Engine; SNIC field segments (0.3--10\,ha)
\item Seeded random sample: \textbf{240 fields} (20 rainfed + 10 irrigated per site)
\item Crop condition: $\mathrm{SVI} = (\mathrm{NDVI} - \mu_{\mathrm{doy}}) / \sigma_{\mathrm{doy}}$ per field, normals to 2023 with the scored year removed
\end{itemize}
\end{card}
\column{0.51\textwidth}
\begin{keycard}[equal height group=sat]{Results}
\small
\renewcommand{\arraystretch}{1.2}
\begin{tabularx}{\linewidth}{@{}Ycc@{}}
\toprule
\textbf{Lag (SVI vs SPEI)} & \textbf{2019--23} & \textbf{2024--26} \\ \midrule
Rainfed: SPEI-30, 10 days & $\rho$ 0.48 & 0.45 \\
Irrigated: SPEI-60, 15 days & 0.24 & 0.29 \\ \midrule
\textbf{15-day forecast} & \textbf{RF} & \textbf{Persist.} \\ \midrule
Macro-F1 (SVI class) & \textbf{0.68} & 0.65 \\
Onset AUC (logistic 0.72) & \textbf{0.76} & 0.50 \\
\bottomrule
\end{tabularx}\\[4pt]
{\raggedright Crops respond about \textbf{10 days} after weather drought.\par}
\end{keycard}
\end{columns}
\end{frame}

% ============================================================ 12 REPLAY
\begin{frame}{5.8 Replay against Government Drought Orders}
\begin{columns}[T,onlytextwidth]
\column{0.47\textwidth}
\begin{card}{Rules (no retraining, real data delays)}
\scriptsize
\begin{itemize}
\item Issued every 5 days, June--November, 2023--2026
\item Weather 7 days late, satellite 2 days late
\item \textbf{Alarm}: SPEI-30 $< -1$ or severe 15-day forecast
\item \textbf{Confirmed}: median rainfed SVI $< -1$ ($\ge 5$ clear fields)
\end{itemize}
\end{card}
\begin{keycard}{Results}
\scriptsize
\renewcommand{\arraystretch}{1.05}
\begin{tabular*}{\linewidth}{@{}l@{\extracolsep{\fill}}ccc@{}}
\toprule
\textbf{Official} & \textbf{Confirmed} & \textbf{Before order} & \textbf{Mandal} \\ \midrule
Severe & \textbf{9 / 9} & 5 / 9 & 6 / 6 \\
Moderate & 3 / 6 & 3 / 6 & 2 / 4 \\
Not declared & 1 / 11 (false alarm) & -- & 1 / 8 \\
\bottomrule
\end{tabular*}\\[2pt]
Two severe droughts confirmed \textbf{22 days before} the order (2026). 2025 moderate orders missed: rain was above normal.
\end{keycard}
\column{0.51\textwidth}
\centering
\includegraphics[width=\linewidth]{figures/res_s_cases.png}\\
{\scriptsize SPEI-30 (blue), rainfed SVI (green), alarms (orange), confirmed (red), G.O. date (black).}
\end{columns}
\end{frame}

% ============================================================ 13 WHY RABI
\begin{frame}{Why a Separate Rabi Module?}
\begin{columns}[T,onlytextwidth]
\column{0.49\textwidth}
\begin{warncard}{1. Different water source}
\footnotesize Rabi crops (Oct--Apr) grow on soil water stored after the monsoon, irrigation and little rain, and often fail late. In-season rain is unrelated to rabi yields ($\rho = +0.01$).
\end{warncard}
\begin{warncard}{2. The kharif chain failed in rabi}
\footnotesize Applied unchanged to rabi 2023--2026, it confirmed \textbf{none of the 4} severe rabi mandal-seasons declared by the government.
\end{warncard}
\column{0.49\textwidth}
\begin{warncard}{3. Too few rabi declarations}
\footnotesize Only 4 severe rabi mandal-seasons in our area: far too few to train or test a model on government orders.
\end{warncard}
\begin{warncard}{4. Harvest looks like drought}
\footnotesize Greenness falls at harvest (Jan--Mar) as well as in drought, so a greenness drop alone cannot be trusted.
\end{warncard}
\end{columns}
\vspace{3mm}
\begin{keycard}{Our solution}
\footnotesize Measure \textbf{root-zone soil moisture}, use it to tell drought from harvest, and validate against an independent outcome: \textbf{district crop yields} (1997--2022).
\end{keycard}
\end{frame}

% ============================================================ 14 LIT RABI
\begin{frame}{6.1 Literature Review -- Rabi Drought}
\footnotesize
\renewcommand{\arraystretch}{1.22}
\begin{tabularx}{\textwidth}{|>{\raggedright\arraybackslash}p{3.3cm}|Y|Y|Y|}
\hline
\textbf{Paper} & \textbf{Contribution} & \textbf{Limitation} & \textbf{Our response} \\ \hline
\textbf{EDI + remote sensing} \newline Padhee et al., 2014 & Rabi drought in Bundelkhand from rain, soil moisture and NDVI lags & Linear, regional; not validated against crop outcomes & Soil indicator validated against district crop yields \\ \hline
\textbf{Early-season drought} \newline Nikam et al., 2020 & Early agricultural drought assessment from remote sensing & Index-based only & Independent yield outcomes, early warning by 15 Dec \\ \hline
\textbf{Sentinel-2 crop failure} \newline Descals et al., 2025 & Sentinel-2 indices decline in drought-hit winter cereals & Temperate crops, field survey; harvest not separated & Soil check separates drought from harvest \\ \hline
\textbf{Terminal drought} \newline Pang et al., 2017 & Chickpea water use and yield under terminal (late) drought & Plant-level study & Season-long soil deficit duration as indicator \\ \hline
\end{tabularx}
\begin{card}{Our direction}
\footnotesize Root-zone soil moisture deficit duration + MODIS greenness $\rightarrow$ crop-loss early warning; field early-drying model checked against soil.
\end{card}
\end{frame}

% ============================================================ 15 RABI METHOD
\begin{frame}{6.2 Rabi Module -- Method}
\vspace{1mm}
\begin{columns}[T,onlytextwidth]
\column{0.49\textwidth}
\begin{card}[equal height group=rabim]{Soil drought indicator}
\normalsize\raggedright
Root zone (cm layers):\\
$\theta = 0.07\,\theta_{0\text{--}7} + 0.21\,\theta_{7\text{--}28} + 0.72\,\theta_{28\text{--}100}$\\[4pt]
7-day mean $\rightarrow$ percentile vs the same point and date ($\pm$15 days) in 1990--2004.\\[4pt]
\textbf{Deficit day}: percentile $\le 0.20$.\\[2pt]
\textbf{DD} = share of 1 Nov -- 28 Feb days in deficit.
\end{card}
\column{0.49\textwidth}
\begin{card}[equal height group=rabim]{Outcome (never used for training)}
\normalsize
\begin{itemize}
\item Multi-crop rabi \textbf{yield index}: each crop's yield vs the median of its previous 5 seasons, fixed area weights
\item \textbf{Loss season}: index $\le -15\%$
\item Early-warning inputs end 8 days before the issue date
\end{itemize}
\end{card}
\end{columns}
\end{frame}

% ============================================================ 16b RABI ARCHITECTURE
\begin{frame}{6.3 Rabi Module -- Architecture}
\centering
\includegraphics[height=0.76\textheight]{figures/v2_rabi_slide.png}
\end{frame}

% ============================================================ 16 RABI RESULTS
\begin{frame}{6.4 Rabi Module -- Results against Crop Yields}
\begin{columns}[T,onlytextwidth]
\column{0.5\textwidth}
\scriptsize
\renewcommand{\arraystretch}{1.25}
\begin{tabular}{@{}lcc@{}}
\toprule
\textbf{Indicator} (Kurnool + Anantapur) & \textbf{$\rho$ with yield} & \textbf{Loss AUC, 15 Dec} \\ \midrule
\textbf{Soil deficit duration (DD)} & $\mathbf{-0.58}$ & \textbf{0.80} \\
Soil moisture at sowing & $+0.55$ & 0.77 \\
Rain June--November & $+0.40$ & 0.72 \\
SPEI-30 & $+0.44$ & 0.60 \\
Rain Nov--Feb (in season) & $+0.01$ & -- \\
\bottomrule
\end{tabular}\\[2pt]
{\scriptsize 46 district-seasons 2000--2022, 11 loss seasons; 95\% CI by resampling seasons.}
\begin{keycard}{Key findings}
\footnotesize
\begin{itemize}
\item Soil deficit warns of crop-loss seasons by \textbf{15 December} (AUC 0.80) --- SPEI only 0.60.
\item \textbf{MODIS greenness} confirms on 13 independent seasons (2010--22): AUC \textbf{0.83}, $\rho$ = +0.70.
\item Two losses with wet soil and green crops are flagged by neither: losses not caused by drought.
\end{itemize}
\end{keycard}
\column{0.48\textwidth}
\centering
\includegraphics[width=\linewidth]{figures/rabi_auc.png}\\[2pt]
\begin{warncard}{Limitation}
\scriptsize Districts became greener after 2010 (Dec--Feb NDVI 0.32 $\rightarrow$ 0.40), so a fixed 2000--09 greenness reference rarely flags recent seasons.
\end{warncard}
\end{columns}
\end{frame}

% ============================================================ 17 FIELD EARLY DRYING
\begin{frame}{6.5 Rabi -- Field Early-Drying Warnings with a Soil Check}
\begin{columns}[T,onlytextwidth]
\column{0.44\textwidth}
\begin{card}{Model}
\footnotesize
\begin{itemize}
\item 1,350 random Sentinel-2 cropland points
\item \textbf{Event}: established crop dries $\ge 21$ days before the point's normal drying date
\item Random forest, 39 causal features (NDVI history, phenology, weather, root-zone soil)
\item Warns 7 / 15 / 30 days ahead; model and thresholds chosen on validation 2022 only
\item Each warning tagged \textbf{soil-dry} or \textbf{soil-wet} using the soil percentile
\end{itemize}
\end{card}
\column{0.54\textwidth}
\scriptsize
\renewcommand{\arraystretch}{1.25}
\begin{tabular}{@{}lccc@{}}
\toprule
\textbf{30 days ahead, test 2023--25} & \textbf{Precision} & \textbf{Recall} & \textbf{Lead} \\ \midrule
Validation-selected model & 0.54 & 0.28 & 16 d \\ \midrule
\multicolumn{4}{@{}l}{\textit{Same events caught (40 of 108):}} \\
Earlier model (no soil) & 0.51 & 0.37 & -- \\
Final model (39 features) & 0.56 & 0.37 & -- \\
\textbf{Final model, soil-dry warnings} & \textbf{0.69} & 0.37 & -- \\
\bottomrule
\end{tabular}
\begin{keycard}{Key finding}
\footnotesize Warnings under below-normal soil are correct \textbf{about twice as often} as under wet soil (0.67 vs 0.38), where greenness drops are mostly harvest.
\end{keycard}
\begin{warncard}{Honest note}
\scriptsize Only 2023 had below-normal soil among the test seasons, so this contrast is descriptive. A separate 15-day NDVI forecast has 18\% lower error than the best baseline.
\end{warncard}
\end{columns}
\end{frame}

% ============================================================ 18 LIT UAV
\begin{frame}{7.1 Literature Review -- UAV Crop Stress Segmentation}
\footnotesize
\renewcommand{\arraystretch}{1.22}
\begin{tabularx}{\textwidth}{|>{\raggedright\arraybackslash}p{3.3cm}|Y|Y|Y|}
\hline
\textbf{Paper} & \textbf{Contribution} & \textbf{Limitation} & \textbf{Our response} \\ \hline
\textbf{SegFormer} \newline Xie et al., 2021 & Hierarchical Transformer encoder with lightweight MLP decoder & General-purpose, 3-channel RGB & 9 channels (RGB + vegetation indices), 7 crop classes \\ \hline
\textbf{Cotton water stress from UAV} \newline Shahi et al., 2023 & CNN classifies cotton water stress from UAV RGB images & Image-level classes; no area & Pixel-level segmentation and area in m$^2$ \\ \hline
\textbf{UAV weed detection} \newline Dang et al., 2023 & Compares deep segmentation networks for weeds in UAV images & Dry weeds and drought look alike & Separate dry-weed and drought classes \\ \hline
\textbf{Lov\'asz-Softmax / focal loss} \newline Berman 2018; Lin 2017 & Losses for IoU and class imbalance & General methods; each alone fixes one problem (imbalance or IoU), not rare drought patches & Focal + Dice + Lov\'asz combined; drought class weight 10 \\ \hline
\end{tabularx}
\begin{card}{Our direction}
\footnotesize 9-channel UAV orthomosaics $\rightarrow$ SegFormer-B2 $\rightarrow$ 7-class map $\rightarrow$ seam-free stitching of billion-pixel fields $\rightarrow$ area and visual stress flag.
\end{card}
\end{frame}

% ============================================================ 19 UAV DATA
\begin{frame}{7.2 UAV Data Acquisition and Photo Stitching}
\vspace{1mm}
\begin{columns}[T,onlytextwidth]
\column{0.36\textwidth}
\begin{card}[equal height group=uavdata]{Flight plan}
\footnotesize
\begin{itemize}
\item DJI Mini 4 Pro (RGB camera)
\item Altitude 10--12\,m, speed 1.2--1.5\,m/s
\item Gimbal $-90^\circ$, path spacing 3--5\,m
\item Overlap 70--80\%, east--west lines
\item Flights: Nov 2025 -- Jan 2026
\end{itemize}
\vspace{1mm}
\centering
\includegraphics[height=1.9cm]{figures/prof/lawnmower_path.png}\\[1pt]
{\tiny\color{muted} Lawnmower path (source: Humboldt State Univ.)}
\end{card}
\column{0.62\textwidth}
\begin{card}[equal height group=uavdata]{Photos $\rightarrow$ orthomosaic (Agisoft Metashape)}
\centering
\begin{minipage}[c]{2.8cm}\centering
\includegraphics[width=\linewidth]{figures/prof/p11_0.png}\\[2pt]
\includegraphics[width=\linewidth]{figures/prof/p11_2.png}
\end{minipage}\hspace{3mm}
{\LARGE$\Rightarrow$}\hspace{3mm}
\begin{minipage}[c]{1.9cm}\centering
\includegraphics[height=3.2cm]{figures/prof/ortho_cotton.png}
\end{minipage}\\[5pt]
\raggedright\small Overlapping photos are matched and aligned into one orthomosaic.\\[2pt]
\textbf{14 orthomosaics}: 3.8 million -- 2.86 billion pixels; GSD 1.68--5.55\,mm (9 fields), 3.6--18.5\,mm (5 Peddakadubur fields, low-resolution frames).
\end{card}
\end{columns}
\end{frame}

% ============================================================ 20 PIPELINE
\begin{frame}{7.3 UAV Module -- End-to-End Pipeline}
\centering
\includegraphics[height=0.76\textheight]{figures/v2_uav_pipeline_slide.png}
\end{frame}

% ============================================================ 21 LABELS + CHANNELS
\begin{frame}{7.4 Building Labels and the Nine Input Channels}
\centering
\includegraphics[width=0.80\textwidth]{figures/uav_step_labels.png}\\[0.5mm]
\begin{columns}[T,onlytextwidth]
\column{0.45\textwidth}
\begin{card}{Labels}
\scriptsize
\begin{itemize}
\item Pseudo-labels: 512\,px window, stride 256, dense CRF; confidence $< 0.55$ $\rightarrow$ ignore
\item Custom tool (Tkinter, Matplotlib, OpenCV): brush 2--150\,px, 50-step undo
\item \textbf{279 clean tiles} (224 train, 55 validation)
\item Per epoch: 500 clean + 500 pseudo + 1,000 synthetic (weights 8 / 0.5 / 1)
\end{itemize}
\end{card}
\column{0.53\textwidth}
\begin{card}{Nine channels (RGB in [0,1])}
\scriptsize
RGB (ImageNet-normalised) \;+\; $\tfrac{2G-R-B}{2}$ \;+\; $\tfrac{(2G-R-B)-(1.4R-G)}{3}$ \;+\; $-\tfrac{0.441R-0.811G+0.385B}{0.811}$\\[2pt]
+\; $\mathrm{GLI}=\tfrac{2G-R-B}{2G+R+B}$ \;+\; $\mathrm{NGRDI}=\tfrac{G-R}{G+R}$ \;+\; $V' = 2V-1$\\[2pt]
\includegraphics[width=\linewidth]{figures/uav_step_channels.png}
\end{card}
\end{columns}
\end{frame}

% ============================================================ 22 MODEL + TRAINING
\begin{frame}{7.5 Model and Training -- Nine-Channel SegFormer-B2}
\vspace{1mm}
\begin{columns}[c,onlytextwidth]
\column{0.61\textwidth}
\centering
\includegraphics[width=\linewidth]{figures/v2_uav_model_slide.png}
\column{0.37\textwidth}
\begin{card}{Training}
\scriptsize\raggedright
$\mathcal{L} = 1.0\,\mathcal{L}_{\text{focal}}(\gamma{=}2) + 2.5\,\mathcal{L}_{\text{Dice}}$\\
$\phantom{\mathcal{L}} \quad + 1.0\,\mathcal{L}_{\text{Lov\'asz}}$\\[3pt]
$\mathcal{L}_{\text{focal}} = -\alpha_y (1-p_y)^{\gamma}\log p_y$\\[3pt]
Class weights $\alpha_c = N/(C N_c)$, clipped to [0.1, 50]; drought 10, cotton boll 8.\\[3pt]
AdamW, lr $2\times10^{-6}$, cosine decay, batch 16, 10 epochs; best validation mIoU at epoch 5.
\end{card}
\begin{keycard}{Logit adjustment}
\scriptsize\raggedright Drought $-1.2$, dry weed $-0.4$ (chosen on validation tiles):\\[2pt]
drought precision \textbf{54.4\% $\rightarrow$ 76.4\%}\\
recall 95.5\% $\rightarrow$ 82.1\%\\
drought IoU \textbf{53.0\% $\rightarrow$ 65.5\%}
\end{keycard}
\end{columns}
\end{frame}

% ============================================================ 23 STITCHING
\begin{frame}{7.6 Prediction Stitching, Area and Stress Decision}
\centering
\includegraphics[width=0.85\textwidth]{figures/uav_step_stitching.png}\\[0.5mm]
\begin{columns}[T,onlytextwidth]
\column{0.5\textwidth}
\begin{card}{Seam-free, memory-safe stitching}
\scriptsize
$1024$\,px tiles, stride 768: \; $A_c(x) = \sum_i w_i(x)\,p_{i,c}(x)$, class $=\arg\max_c A_c(x)$\\
Rolling band: \textbf{1.5\,GB instead of 80\,GB} for a 2.86-billion-pixel orthomosaic.\\
Post-processing: no-data border, glare/shadow fix, small regions, harvest zones.
\end{card}
\column{0.48\textwidth}
\begin{card}{Area and visual stress flag}
\scriptsize
Area $= N_c \times (\mathrm{GSD}/1000)^2$ m$^2$.\\
$D$ = drought share of crop area, $N$ = same outside harvest zones, $C$ = share in the largest drought region.\\
\textbf{Flag} if $D > 0.15$ or ($N > 0.10$ and $C > 0.25$).
\end{card}
\end{columns}
\end{frame}

% ============================================================ 24 VALIDATION
\begin{frame}{7.7 UAV Validation Results (55 validation tiles)}
\begin{columns}[c,onlytextwidth]
\column{0.50\textwidth}
\begin{card}{Model versions (validation mIoU)}
\footnotesize
\renewcommand{\arraystretch}{1.2}
\begin{tabular*}{\linewidth}{@{}l@{\extracolsep{\fill}}r@{}}
\toprule
\textbf{Model version} & \textbf{mIoU (\%)} \\ \midrule
V8, SegFormer-B2 & 61.18 \\
V9 / V9.1, SegFormer-B3 & 44.25 / 45.89 \\
V10, DeepLabV3+ & 44.79 \\
\textbf{V11, SegFormer-B2 (final)} & \textbf{61.46} \\
V12, SegFormer-B2 & 61.39 \\
\bottomrule
\end{tabular*}
\end{card}
\begin{keycard}{Final V11 (after logit adjustment)}
\centering
\begin{tabular*}{\linewidth}{@{\extracolsep{\fill}}ccc@{}}
{\color{navy}\bfseries\small 65.1\%} & {\color{navy}\bfseries\small 82.7\%} & {\color{navy}\bfseries\small 79.2\%} \\
{\scriptsize mIoU} & {\scriptsize pixel accuracy} & {\scriptsize drought F1}
\end{tabular*}\\[4pt]\raggedright
\scriptsize Per-class IoU: soil 74, crop 73, cotton 53, green weed 71, dry weed 54, drought 65\%.
\end{keycard}
\column{0.48\textwidth}
\centering
\includegraphics[height=0.68\textheight]{figures/uav_confusion.png}\\[1pt]
{\scriptsize\color{muted} Row-normalised confusion matrix (\%)}
\end{columns}
\end{frame}

% ============================================================ 25 FIELD RESULTS
\begin{frame}{7.8 UAV Results on the 14 Fields}
\begin{columns}[T,onlytextwidth]
\column{0.52\textwidth}
\scriptsize
\renewcommand{\arraystretch}{1.15}
\begin{tabular}{@{}lrrrc@{}}
\toprule
\textbf{Field} & \textbf{GSD mm} & \textbf{Area m$^2$} & \textbf{Drought \%} & \textbf{Flag} \\ \midrule
Lepakshi & 2.08 & 2,976 & \textbf{71.53} & \textbf{yes} \\
Sira & 2.14 & 4,341 & \textbf{54.26} & \textbf{yes} \\
Somandepalle & 2.35 & 3,936 & \textbf{25.45} & \textbf{yes} \\
KNL Puddur & 1.73 & 3,866 & 0.15 & no \\
KNL Maize Tower & 1.68 & 2,241 & 0.07 & no \\
Pulliah & 2.69 & 2,897 & 0.04 & no \\
Balaji Villas & 1.84 & 8,044 & 0.06 & no \\
Shilpa & 1.94 & 6,395 & 0.00 & no \\
KNL DDP & 5.55 & 19,812 & 0.00 & no \\
Peddakadubur F1/F3/F4/F6/F7 & 3.6--18.5 & 42--660 & 0--0.89 & no \\
\bottomrule
\end{tabular}\\[2pt]
\begin{keycard}{Key finding}
\scriptsize The 3 flagged fields are exactly the southern fields the survey team judged drought-hit; no Kurnool field is flagged. Largest orthomosaic (2.86 billion px): 45 min on one L4 GPU.
\end{keycard}
\column{0.48\textwidth}
\centering
\includegraphics[height=2.75cm]{figures/uav_lepakshi_overlay.jpg}\\[1mm]
\includegraphics[height=2.55cm]{figures/uav_sira_overlay.jpg}\hspace{2mm}
\includegraphics[height=2.55cm]{figures/uav_knlddp_overlay.jpg}\\
{\scriptsize\color{muted} Lepakshi, Sira (flagged), KNL DDP (not flagged); red = drought}
\end{columns}
\end{frame}

% ============================================================ 26 CONCLUSION
\begin{frame}{8. Conclusion, Limitations and Future Work}
\vspace{1mm}
\begin{columns}[T,onlytextwidth]
\column{0.32\textwidth}
\begin{keycard}[equal height group=conc1]{Kharif}
\centering{\color{navy}\bfseries\normalsize 9 / 9\par}\vspace{1pt}
{\scriptsize\color{muted} severe droughts confirmed, 5 before the order\par}\vspace{4pt}
\raggedright\scriptsize Hybrid weather forecast 0.819 / 0.651 at 7 / 15 days; crops respond about 10 days later.
\end{keycard}
\column{0.32\textwidth}
\begin{keycard}[equal height group=conc1]{Rabi}
\centering{\color{navy}\bfseries\normalsize AUC 0.80\par}\vspace{1pt}
{\scriptsize\color{muted} crop-loss warning by 15 December\par}\vspace{4pt}
\raggedright\scriptsize Confirmed by MODIS (0.83); soil-dry field warnings twice as reliable.
\end{keycard}
\column{0.32\textwidth}
\begin{keycard}[equal height group=conc1]{UAV}
\centering{\color{navy}\bfseries\normalsize 65.1\% mIoU\par}\vspace{1pt}
{\scriptsize\color{muted} 7-class maps of billion-pixel fields\par}\vspace{4pt}
\raggedright\scriptsize Flags the 3 fields the survey team judged drought-hit; 1.5\,GB instead of 80\,GB.
\end{keycard}
\end{columns}
\vspace{3mm}
\begin{columns}[T,onlytextwidth]
\column{0.49\textwidth}
\begin{warncard}[equal height group=conc2]{Limitations}
\scriptsize
\begin{itemize}
\item No field-level drought ground truth; orders and yields are indirect
\item No skill 30 days ahead (kharif); few dry rabi test seasons
\item UAV: validation tiles only; RGB shows visual stress, not its cause
\end{itemize}
\end{warncard}
\column{0.49\textwidth}
\begin{card}[equal height group=conc2]{Future work}
\scriptsize
\begin{itemize}
\item Run the full kharif / rabi $\rightarrow$ UAV loop in rabi 2026-27
\item Finish the 127-tile independent UAV test with field confirmation
\item Trend-robust greenness reference; multispectral / thermal drone imagery
\end{itemize}
\end{card}
\end{columns}
\end{frame}

% ============================================================ 27 REFERENCES
\begin{frame}{References}
\setbeamerfont{itemize/enumerate body}{size=\fontsize{6.5}{7.7}\selectfont}
\setbeamerfont{enumerate item}{size=\fontsize{6.5}{7.7}\selectfont}
\fontsize{6.5}{7.7}\selectfont
\begin{columns}[T,onlytextwidth]
\column{0.49\textwidth}
\begin{enumerate}\setlength{\itemsep}{0.6pt}
\item Govt.\ of India, \emph{Manual for Drought Management}, 2016 (updated 2020).
\item K.~Bhardwaj, V.~Mishra, ``Drought detection and declaration in India,'' \emph{Water Security}, 14, 2021.
\item S.~M.~Vicente-Serrano \emph{et al.}, ``A multiscalar drought index sensitive to global warming: the SPEI,'' \emph{J.\ Climate}, 23(7), 2010.
\item A.~Dikshit, B.~Pradhan, A.~Huete, ``An improved SPEI drought forecasting approach using the LSTM,'' \emph{J.\ Environ.\ Manage.}, 283, 2021.
\item A.~Marusov \emph{et al.}, ``Long-term drought prediction using deep neural networks based on geospatial weather data,'' \emph{Environ.\ Model.\ Softw.}, 179, 2024.
\item A.~Belayneh \emph{et al.}, ``Long-term SPI drought forecasting \ldots\ wavelet neural network and wavelet SVR,'' \emph{J.\ Hydrology}, 508, 2014.
\item H.~Boogaard \emph{et al.}, ``AgERA5,'' Copernicus Climate Data Store.
\item F.~N.~Kogan, ``Application of vegetation index and brightness temperature for drought detection,'' \emph{Adv.\ Space Res.}, 15(11), 1995.
\item A.~J.~Peters \emph{et al.}, ``Drought monitoring with NDVI-based Standardized Vegetation Index,'' \emph{PE\&RS}, 68(1), 2002.
\item S.~Park \emph{et al.}, ``Drought assessment and monitoring through blending of multi-sensor indices using machine learning,'' \emph{Agric.\ Forest Meteorol.}, 216, 2016.
\item S.~S.~Sharma \emph{et al.}, ``Leveraging Sentinel-2 data and machine learning for drought detection in India,'' \emph{Remote Sens.}, 17(18), 2025.
\end{enumerate}
\column{0.49\textwidth}
\begin{enumerate}\setcounter{enumi}{11}\setlength{\itemsep}{0.6pt}
\item N.~Gorelick \emph{et al.}, ``Google Earth Engine,'' \emph{Remote Sens.\ Environ.}, 202, 2017.
\item S.~K.~Padhee \emph{et al.}, ``Integrating EDI and remote sensing derived parameters for agricultural drought assessment in Bundelkhand,'' \emph{ISPRS Arch.}, XL-8, 2014.
\item B.~R.~Nikam \emph{et al.}, ``Assessment of early season agricultural drought using remote sensing,'' \emph{ISPRS Arch.}, XLIII-B3, 2020.
\item A.~Descals \emph{et al.}, ``Evaluating Sentinel-2 for monitoring drought-induced crop failure in winter cereals,'' \emph{Remote Sens.}, 17(2), 2025.
\item J.~Pang \emph{et al.}, ``Pattern of water use and seed yield under terminal drought in chickpea,'' \emph{Front.\ Plant Sci.}, 8, 2017.
\item E.~Xie \emph{et al.}, ``SegFormer: simple and efficient design for semantic segmentation with Transformers,'' \emph{NeurIPS}, 2021.
\item T.~B.~Shahi \emph{et al.}, ``Classification of water stress in cotton using CNNs from UAV images,'' \emph{Remote Sens.}, 15(1), 2023.
\item F.~Dang \emph{et al.}, ``Deep learning-based weed detection using UAV images: a comparative study,'' \emph{Drones}, 7(1), 2023.
\item M.~Berman \emph{et al.}, ``The Lov\'asz-Softmax loss,'' \emph{CVPR}, 2018.
\item T.-Y.~Lin \emph{et al.}, ``Focal loss for dense object detection,'' \emph{ICCV}, 2017.
\item R.~Achanta, S.~S\"usstrunk, ``Superpixels and polygons using simple non-iterative clustering,'' \emph{CVPR}, 2017.
\end{enumerate}
\end{columns}
\end{frame}

% ============================================================ 28 THANK YOU
\begin{frame}[plain]
\centering
\vspace{8mm}
\includegraphics[height=2.0cm]{figures/logo.jpeg}\\[5mm]
{\fontsize{34}{38}\selectfont\color{navy} Thank You}\\[5mm]
{\large A Multimodal Framework for Agricultural Drought Monitoring and Prediction}\\[3mm]
{\normalsize Weather Time Series \;$|$\; Satellite Imagery \;$|$\; UAV Semantic Segmentation}\\[5mm]
{\footnotesize Department of Computer Science and Engineering\\
Indian Institute of Information Technology Design and Manufacturing Kurnool}
\end{frame}

\end{document}
